\documentclass[10pt,a4paper]{amsart}
\usepackage[margin=19mm]{geometry}
\usepackage{amsmath,amssymb,mathtools}
\usepackage[scr=rsfs,cal=euler]{mathalfa}
\usepackage{xfrac}
\usepackage[unicode,hidelinks,bookmarks=false]{hyperref}
\newcommand{\ie}{i.e.\ }
\newcommand{\cf}{cf.\ }
\newcommand{\resp}{respectively\ }
\newcommand{\Subsection}{\subsection}
\providecommand{\nobreakdash}{\nobreak\hbox{-}}
\setlength{\emergencystretch}{2em}
\multlinegap0pt
\usepackage{bm}
\usepackage{bbm}
\usepackage{dsfont}
\usepackage{xspace}
\usepackage{cancel}
\usepackage{mathtools}
\usepackage{amsthm}
\makeatletter
\@ifundefined{theorem}{\newtheorem{theorem}{Theorem}}{}
\@ifundefined{proposition}{\newtheorem{proposition}[theorem]{Proposition}}{}
\makeatother
\makeatletter
\newcommand\K{\ensuremath{\mathbb{K}}}
\expandafter\ifx\csname assumption\endcsname\relax
\newenvironment{assumption}[1]
 {\renewcommand\thetaggedtheoremx{#1}\taggedtheoremx}
\fi
\expandafter\ifx\csname theoremp\endcsname\relax
\newenvironment{theoremp}[1]{%
  \renewcommand\thetheoremalttmp{#1}%
  \edef\@currentlabel{#1}%
  \theoremalttmp
}{%
  \endtheoremalttmp
}
\fi
\makeatother
\title[Sparse supports of lattice eigenfunctions: quantitative grow]{Sparse supports of lattice eigenfunctions: quantitative growth and algebraic rigidity}
\author{Yunlei Wang}
\date{Source: arXiv:2608.01673v2 (CC BY 4.0)}
\begin{document}
\maketitle

\noindent\emph{Source and licence.} Sparse supports of lattice eigenfunctions: quantitative growth and algebraic rigidity. Version arXiv:2608.01673v2 (CC BY 4.0), distributed under the Creative Commons Attribution 4.0 International licence, as recorded in the arXiv licence metadata for this exact version and shown on the abstract page https://arxiv.org/abs/2608.01673v2. Full rights notice: licenses/arxiv-2608.01673-CC-BY-4.0.txt.\par\smallskip

\noindent\textit{Editorial scope.} Counted unit: the Proposition whose source label is \texttt{prop:norm-defect} in the article (source0.tex lines 1670--1677, source label \texttt{prop:norm-defect}), delivered together with the article's own complete original proof of it (source0.tex lines 1679--1767, one \texttt{proof} environment of 89 lines). No mathematics is written, repaired or re-derived: statement and proof are the article's own text line for line. Required context retained uncounted: thm:green-boundary (lines 1492--1508). Every definition, display and label of those lines is retained unchanged. Omitted from this package and not redistributed: 1--1491, 1509--1669, 1678--1678, 1768--2715. Nothing inside a retained excerpt is omitted or altered: every retained body is delivered exactly as the article's lines are written, so no omission receipt of any kind accompanies this package. Citations: the retained excerpts carry no citation command at all, so no bibliography entry of the article is transcribed and no reference is redistributed. Reference adaptation: none was needed and none was made; every reference inside the delivered unit resolves inside it, so no unresolved reference or citation is produced. Printed slips kept literal: no printed word, symbol, hypothesis or number of the counted statement or of its proof was altered. Layout and class: the article's own class is \texttt{amsart} with packages including xcolor, hyperref, yhmath, geometry, lmodern, fontenc, textcomp, amsmath, amssymb, mathrsfs, mathtools, tikz, orcidlink, bbm; the wrapper is the lane's pinned \texttt{amsart} editorial wrapper at 10pt on A4, which restates the theorem-like environments the retained text needs and adds the article's own macro definitions that the retained text needs (\texttt{\textbackslash K}), with extra wrapper packages (bm, bbm, dsfont, xspace, cancel, mathtools). The delivered numbering is the unit's own. Assets: no retained span contains a figure environment, a graphics-inclusion command, a PDF-page inclusion, a table or a TikZ picture; no image file is copied into this package, the package declares no asset dependency and no image file is redistributed with it. Licence: version arXiv:2608.01673v2 of this article is distributed under the Creative Commons Attribution 4.0 International licence, as recorded in the arXiv licence metadata for this exact version and shown on the abstract page https://arxiv.org/abs/2608.01673v2. A full rights notice is delivered at licenses/arxiv-2608.01673-CC-BY-4.0.txt. Characters: every retained excerpt is pure ASCII, so the delivered engine, XeLaTeX with its default Latin Modern text font, reports no missing character.
\begin{theorem}\label{thm:green-boundary}
Let $\mathscr{P}$ and $\mathfrak{n}$ be given above,
and let the one-dimensional torus act by
$p_i\mapsto s^{w_i}p_i$, where the positive weights $w_i$ are pairwise
distinct.  Suppose that $M$ is a finite-length equivariant
$\mathscr P$-module generated by at most $r$ elements.  Let $\Gamma\in \mathscr{P}$ be a polynomial with linear part $\ell=\sum_{i=1}^dc_ip_i$ where $c_i\neq 0$, that is,
\[
 \Gamma=\ell+h,\quad \text{for some }h\in \mathfrak{n}^2.
\]
then
\begin{equation}\label{eq:weighted-green-boundary}
 \dim_\K M/\Gamma M
 \le C_{d,r}\bigl(\dim_\K M\bigr)^{(d-1)/d}.
\end{equation}
The statement is unchanged if $\Gamma$ is a formal series, since a power
of $\mathfrak{n}$ annihilates $M$.
\end{theorem}

\begin{proposition}
\label{prop:norm-defect}
For every homogeneous quotient $A=\K[V_1,\ldots,V_d]/J$ and every
$t\ge0$,
\begin{equation}\label{eq:norm-defect}
        k_A(t)\le C_dH_A(t)^{(d-1)/d}.
\end{equation}
\end{proposition}

\begin{proof}
Let
\begin{equation*}
    \mathfrak{m}=(V_1,\ldots,V_d)
\end{equation*}
and set
\begin{equation*}
    A':=A/\mathfrak{m}^{t+1}A,
    \quad
    \mathcal K:=
    \left(0:_{A'[z]}\widetilde{\mathcal R}_d\right).
\end{equation*}
The quotient map $A\to A'$ is an isomorphism through degree $t$.
It therefore does not change the degree-$t$ source of multiplication by
$\widetilde{\mathcal R}_d$, although it may introduce additional kernel
elements. Hence
\begin{equation*}
    k_A(t)
    \leq
    \dim_{\K}\mathcal K_t.
\end{equation*}
The module $A'[z]$ is finite free over $\K[z]$. Its graded submodule
$\mathcal K$ is torsion free and hence free over the principal ideal
domain $\K[z]$. Since its homogeneous generators have nonnegative
degrees,
\begin{equation*}
    \dim_{\K}\mathcal K_t
    \leq
    \operatorname{rank}_{\K[z]}\mathcal K.
\end{equation*}
Consequently,
\begin{equation*}
    k_A(t)
    \leq
    \dim_{\K}\mathcal K_t
    \leq
    \operatorname{rank}_{\K[z]}\mathcal K.
\end{equation*}




Set
\[
 p_i=\sum_{j=1}^dV_j^{2i},\qquad 1\le i\le d.
\]
Consider the signed-permutation group, denoted by $W=(\mathbb Z/2\mathbb Z)^d\rtimes S_d$, which changes signs and
permutes the variables $V_i$. Since $p_1,\ldots,p_d$ are basic invariants for the signed-permutation
group, $\mathbb K[V]$ is free of rank $2^d d!$ over
$\mathbb K[p_1,\ldots,p_d]$. Hence $A'$ is generated by at most
$2^d d!$ elements over this ring.
Scalar dilation acts with the distinct
weights $2,4,\ldots,2d$, so $A'$ satisfies the equivariance and bounded
generator hypotheses of Theorem~\ref{thm:green-boundary}.  In the completed
local ring, all conjugate factors of the norm except the harmonic one are
units, while the remaining factor is
\begin{equation}\label{eq:norm-linearization}
 \sum_{j=1}^d\sqrt{4+V_j^2}-2d
   =\sum_{i=1}^da_ip_i+O((p_1,\ldots,p_d)^2),
 \qquad a_i\ne0
\end{equation}
Here we absorb the higher power sums into the error term by Newton identities.
The homogeneous ideal defining $A'$ is preserved by the invertible change
$V=zW$.  Over $\K(z)$ we therefore have the exact chain
\begin{equation}\label{eq:norm-rank-cokernel}
 \operatorname{rank}_{\K[z]}(0:\widetilde{\mathcal R}_d)
 =\dim_{\K(z)}\ker(\mathcal R_d:A'_{\K(z)}\to A'_{\K(z)})
 =\dim_{\K(z)}A'_{\K(z)}/\mathcal R_dA'_{\K(z)}.
\end{equation}
Since $(V_1,\ldots,V_d)^{t+1}A'=0$, each formal square root
\begin{equation*}
    \omega_i=\sqrt{4+V_i^2}
    =2+\frac14V_i^2-\frac1{64}V_i^4+\cdots
\end{equation*}
truncates to a well-defined element of $A'$. The norm factorization
therefore holds directly in $A'$. Every factor except
\begin{equation*}
    \Gamma=\sum_{i=1}^d\omega_i-2d
\end{equation*}
has nonzero constant term and is consequently a unit. Hence
$\mathcal R_d=U\Gamma$ for some unit $U$, so multiplication by
$\mathcal R_d$ and multiplication by $\Gamma$ have the same kernel and
cokernel. Thus Theorem~\ref{thm:green-boundary}
bounds \eqref{eq:norm-rank-cokernel} by
\begin{equation*}
    C_d(\dim A')^{(d-1)/d}=C_dH_A(t)^{(d-1)/d}.
\end{equation*}
This proves~\eqref{eq:norm-defect}.
\end{proof}

\end{document}
